\subsection{Phase 1: A philosophical shield against magical thinking}\label{03_it_started-phase-1-a-philosophical-shield-against-magical-thinking}

Before I start talking about signaling, secretory cells, early nervous systems, or any of the biological machinery that eventually made minds possible, I should probably say what kind of claim I think I am making here, and what kind I am not.

I am not claiming to have solved the mind. I am not claiming that twentieth- and twenty-first-century philosophy of mind can be cleanly folded into a single manifesto chapter and come out looking tidy. And I am definitely not claiming that the latest computational vocabulary has finally revealed the essence of thought, in contrast to all the previous metaphors that now look dated or embarrassing. What I \emph{am} claiming is more limited, and for that reason maybe more important: if we want to think clearly about agency, failure, and responsibility, then we need concepts that make minds look more like implementations and less like spells.

That does not settle everything. But it rules some things out, and I think ruling them out matters.

\subsubsection{Marr: computation can be described at more than one level}\label{03_it_started-marr-computation-can-be-described-at-more-than-one-level}

David Marr's \emph{Vision} (1982) proposed that information-processing systems can be described at three levels: the \textbf{computational} level, which asks what problem is being solved and why; the \textbf{algorithmic} level, which asks what representations and procedures carry out that solution; and the \textbf{implementational} level, which asks what physical substrate actually realizes the algorithm. His famous line was that the nature of the computation depends more upon the problem being solved than upon the particular hardware used to solve it (\href{https://pillowlab.princeton.edu/pubs/Pillow2024_MarrDebate.pdf}{Pillow Lab, ``The Marr Debate''}).

I find this useful not because it proves anything by itself, but because it makes certain forms of magical thinking harder to sustain. If a capacity can be analyzed at these different levels, then it begins to look less like an occult faculty and more like something that is realized in machinery, even if that machinery is very complicated and even if our present descriptions are crude.

A modern reader is tempted, at this point, to start importing reinforcement learning and prediction-error talk everywhere. Sometimes that temptation is illuminating; sometimes it is just showing off. Later on I \emph{will} come back to that temptation, especially when particular neural systems enter the story. For now, the important point is simpler: Marr gives us a way to talk about mental function without pretending that explanation stops at either pure abstraction or pure tissue. For examples of how later authors map reinforcement learning across Marr's levels, see Niv and Langdon (2016), Schultz, Dayan and Montague (1997), and Dabney et al.~(2020), as cited in the research draft.

\subsubsection{Putnam: function is not identical to one particular substance}\label{03_it_started-putnam-function-is-not-identical-to-one-particular-substance}

Hilary Putnam's 1967 argument against simple mind-brain identity theory made a different but complementary point. Creatures with very different nervous systems can still appear to instantiate similar mental states. If pain, for example, were strictly identical to one very specific physical-chemical state, then it would become mysteriously difficult to explain how organisms with radically different neural architectures could nevertheless count as feeling pain. Putnam's answer was \textbf{multiple realizability}: the same functional state may be realized in different physical substrates (\href{https://plato.stanford.edu/entries/multiple-realizability/}{Stanford Encyclopedia of Philosophy: Multiple Realizability}; \href{https://en.wikipedia.org/wiki/Multiple_realizability}{Wikipedia: Multiple realizability}). Or, in his own memorable formulation, we could in principle be made of Swiss cheese and it would not matter (\href{https://www.britannica.com/biography/Hilary-Putnam/Philosophy-of-mind}{Britannica: Hilary Putnam, philosophy of mind}).

That idea opened the door to \textbf{functionalism}, where what matters is less the material itself than the causal role being played (\href{https://en.wikipedia.org/wiki/Multiple_realizability}{Wikipedia: Multiple realizability}). Jerry Fodor pushed this further by arguing that psychological categories do not line up neatly with the kinds used in lower-level physics, because the same function may be realized in many different ways (see again the \href{https://plato.stanford.edu/entries/multiple-realizability/}{Stanford Encyclopedia of Philosophy entry on Multiple Realizability}). Again, I am not presenting this as settled metaphysical truth. I am treating it as a useful constraint on thought. If mental capacities are at least partly individuated by what they do, then we are already some distance away from older ideas in which character, will, or spirit float free of implementation.

At the same time, multiple realizability cuts both ways. It does \textbf{not} mean implementation is unimportant. On the contrary: if a given function happens, in a human being, to depend on a particular bit of biological machinery, then damage to that machinery matters whether or not the same function might have been realized differently in some alien nervous system, future machine, or philosophical thought experiment.

\subsubsection{The Churchlands: folk psychology is not sacred just because it is familiar}\label{03_it_started-the-churchlands-folk-psychology-is-not-sacred-just-because-it-is-familiar}

Paul Churchland's eliminative materialism is more aggressive, and probably more irritating, but I think it earns its place here. His claim was not just that folk psychology is incomplete, but that some of its basic categories may be so misleading that they will eventually be displaced rather than smoothly translated into mature neuroscience (\href{https://en.wikipedia.org/wiki/Paul_Churchland}{Wikipedia: Paul Churchland}; \href{https://www.sfu.ca/~kathleea/docs/Eliminative\%20materialism.pdf}{Churchland 1981 PDF}). Patricia Churchland's \emph{Neurophilosophy} turned that into a broader program: if philosophy of mind wants to talk seriously about minds, it cannot remain indifferent to what brains are actually like (\href{https://ithy.com/article/patricia-churchland-summary-i5g7z3k2}{Ithy summary}; \href{https://philopedia.org/thinkers/patricia-smith-churchland/}{Philopedia}).

There are at least three reasons this matters. First, folk-psychological language often fails precisely where suffering becomes most serious: it is bad at explaining dreaming, mental illness, creativity, intelligence differences, and many forms of cognitive breakdown (\href{https://philosophy.hku.hk/courses/old/ronmallon/phil2230/phil2230l13.html}{HKU Philosophy notes}). Second, it is strikingly stagnant. Human beings have been talking in roughly the same way about belief, desire, weakness, and will for a very long time (\href{https://plato.stanford.edu/entries/materialism-eliminative/}{Stanford Encyclopedia of Philosophy: Eliminative Materialism}). Third, many of these categories fit awkwardly, if at all, with the rest of what we know from biology, chemistry, and physics (\href{https://en.wikipedia.org/wiki/Eliminative_materialism}{Wikipedia: Eliminative materialism}). They survive partly because they are intuitive, not because they are especially good.

This matters morally as much as philosophically. Terms like \emph{willpower}, \emph{laziness}, or \emph{lack of character} feel descriptive, but they often function as placeholders for ignorance. They are what we say when we are more interested in preserving a moral drama than in asking how a system actually works. That does not mean ordinary language is useless. It means it should lose its monopoly.

Now, there are specific clinical cases in which the folk-psychological picture breaks down so cleanly that it almost becomes embarrassing to keep using it. Abulia, for example, is literally a disorder of diminished initiative, drive, and action (\href{https://www.ncbi.nlm.nih.gov/books/NBK537093/}{NCBI Bookshelf: Abulia}). In Parkinson's disease, computational modeling has shown changes in parameters related to how reward affects willingness to exert effort, and those changes can be modified by neurochemical state (\href{https://pubmed.ncbi.nlm.nih.gov/27335396/}{PubMed: Le Bouc et al.~2016}; \href{https://pmc.ncbi.nlm.nih.gov/articles/PMC6601748/}{PMC related article}). I am mentioning this here only as a preview of the larger point, not because this section is supposed to become a neurotransmitter argument ahead of schedule. The larger point is that what ordinary language calls ``will'' already begins to fracture the moment one looks at actual mechanisms.

\subsubsection{Why I am carrying all three of these with me}\label{03_it_started-why-i-am-carrying-all-three-of-these-with-me}

Taken together, Marr, Putnam, and the Churchlands do not give us a final theory of mind. What they give us is something more modest and, to me, more useful: a way of making magical thinking less comfortable.

Marr says it can be legitimate to describe a capacity at the level of problem, procedure, and implementation rather than treating it as a mysterious gift. Putnam says function is not reducible to one privileged substance, even though it must still be realized somehow. The Churchlands say our inherited vocabulary for mental life is not holy, and may in fact be actively distorting our understanding. None of that proves a complete picture of agency. But it does support a picture in which agency is \textbf{material, describable, historically evolved, and vulnerable to breakdown}.

That is the only point I need from this section before moving into biology proper. If minds are implemented rather than enchanted, then it makes sense to ask where the relevant machinery came from, what earlier forms it emerged from, what constraints shaped it, and why modern computational analogies---used cautiously, and with full awareness that they too will one day look dated---might still help us think a little more clearly than the old moral language of character and sin. The rest of this section is an attempt to follow that machinery backward.

\subsection{Phase 2: Why computational analogies are tempting, and why I am using them anyway}\label{03_it_started-phase-2-why-computational-analogies-are-tempting-and-why-i-am-using-them-anyway}

I should probably be a little careful here.

This is the point in the story where a person can very easily start sounding as though they have discovered that the universe is secretly doing machine learning. I have not. At most, I think there are places where modern computational language helps illuminate old biological problems a little more clearly than older metaphors did. That does not make the analogy exact, and it definitely does not mean that a cortex is just a stack of tensors with delusions of grandeur. It means only that if evolution and engineering are both forced to solve information-processing problems under severe constraints, then some of the same design patterns may show up in both places.

I am not treating that as revelation. I am treating it as a useful suspicion.

\subsubsection{Why CNNs and visual cortex look uncomfortably alike}\label{03_it_started-why-cnns-and-visual-cortex-look-uncomfortably-alike}

One reason modern machine-learning analogies are hard to resist is that some of them really do seem to recover structures that biologists had already found. Hubel and Wiesel's work on visual cortex helped establish a hierarchical picture in which early visual areas detect relatively simple features such as oriented edges through small receptive fields, while later areas integrate those features into more complex representations (\href{https://en.wikipedia.org/wiki/Convolutional_neural_network}{Wikipedia: Convolutional neural network}; \href{https://people.csail.mit.edu/pulkitag/data/cnn_mimics_brain.pdf}{MIT CSAIL overview}). Convolutional neural networks, developed for engineering rather than neurobiology, ended up relying on a surprisingly similar logic: local receptive fields, hierarchical feature extraction, shared weights, and pooling operations that support a degree of position invariance (\href{https://arxiv.org/pdf/1901.06032}{arXiv review}; \href{https://arxiv.org/html/2412.06740v1}{arXiv overview}).

That parallel does not prove that brains are literally CNNs, and I do not want to write as though it does. What it seems to show, at minimum, is that when two very different processes are trying to solve a similar problem --- extracting robust structure from high-dimensional visual input --- they may converge on comparable strategies. Yamins et al.~(2014) made this especially hard to ignore by showing that a hierarchical CNN optimized only for image categorization, with no explicit constraint to match neural data, nevertheless developed internal representations that predicted responses in inferotemporal cortex and V4 (\href{https://pmc.ncbi.nlm.nih.gov/articles/PMC5780340/}{PubMed Central}). Yamins and DiCarlo later sharpened the point: in this domain, better task performance tended to produce more brain-like internal representations (\href{https://www.researchgate.net/publication/295833538_Using_goal-driven_deep_learning_models_to_understand_sensory_cortex}{ResearchGate summary}).

I do not think this is magic, and I do not think it is coincidence either. It looks more like what happens when different systems are pushed by similar constraints toward workable solutions. Both biological and artificial vision systems have to deal with noisy, high-dimensional inputs, finite resources, and the need to extract invariant structure from changing scenes. Under those conditions, some degree of convergence should perhaps be less shocking than it first sounds. Even so, I would rather say that these results \emph{support} the usefulness of the analogy than that they settle anything metaphysical.

\subsubsection{Compression, efficient coding, and the sheer cost of being a nervous system}\label{03_it_started-compression-efficient-coding-and-the-sheer-cost-of-being-a-nervous-system}

Another reason computational analogies keep showing up is that nervous systems are, among other things, stuck with information bottlenecks. Claude Shannon's work on information theory gave people a formal language for talking about compression, redundancy, and channel capacity. Later theorists such as Horace Barlow proposed that sensory systems might be shaped by pressures toward efficient coding: in other words, toward using limited signaling capacity in ways that preserve as much useful information as possible while wasting as little as possible (\href{https://en.wikipedia.org/wiki/Efficient_coding_hypothesis}{Wikipedia: Efficient coding hypothesis}; \href{https://www.sciencedirect.com/science/article/pii/S0010027716300750}{ScienceDirect overview}).

That idea is not just abstract. The visual system really does compress the world brutally: on the order of \textasciitilde10\^{}8 photoreceptors feeding into \textasciitilde10\^{}6 optic nerve fibers, with only a tiny trickle of that eventually reaching conscious awareness (\href{https://en.wikipedia.org/wiki/Efficient_coding_hypothesis}{Wikipedia: Efficient coding hypothesis}). Laughlin's work on insect photoreceptors and Olshausen and Field's work on sparse coding both helped make the broader point that optimization under natural-scene statistics can recover receptive-field structures that look a great deal like what biology actually built (\href{https://www.mathematicalpsychology.com/Efficient_Coding_Hypothesis}{Mathematical Psychology overview}; \href{https://pubmed.ncbi.nlm.nih.gov/8637596/}{PubMed: Olshausen and Field 1996}; \href{https://scispace.com/papers/emergence-of-simple-cell-receptive-field-properties-by-bulmogedmu}{SciSpace summary}).

The reason this matters for the present chapter is not that I want to reduce perception to an equation. It is that nervous systems are expensive, and the expense is not optional. Spikes cost energy. Synaptic transmission costs energy. The brain uses a wildly disproportionate share of the body's energy budget, and that fact alone should make us suspicious of any picture of the mind that floats serenely above metabolism (\href{https://www.mathematicalpsychology.com/Efficient_Coding_Hypothesis}{Mathematical Psychology overview}; \href{https://www.lesswrong.com/posts/mW7pzgthMgFu9BiFX/the-brain-is-not-close-to-thermodynamic-limits-on}{LessWrong discussion of thermodynamic limits}; \href{https://aleph.se/andart2/neuroscience/energetics-of-the-brain-and-ai/}{Aleph essay on brain energetics}). Some of the stronger formulations in this literature probably deserve to be handled with care, but the general lesson seems hard to escape: whatever else minds are, they are not exempt from bandwidth limits, energetic tradeoffs, and the pressure to encode efficiently.

That, at least, is one reason machine-learning analogies do not strike me as pure gimmickry. They give us a vocabulary for talking about constrained information processing without pretending that matter is incidental.

\subsubsection{Plants, metabolism, and the embarrassment of calling neurotransmitters ``neural''}\label{03_it_started-plants-metabolism-and-the-embarrassment-of-calling-neurotransmitters-neural}

If the previous examples show why computational analogies can be helpful, plants help with a different but related task: they make certain kinds of neural mysticism look faintly ridiculous.

Toyota et al.~(2018) showed that when \emph{Arabidopsis} leaves are wounded, glutamate released from damaged cells can trigger calcium waves that propagate through the plant and help coordinate defense responses in distant leaves (\href{https://news.wisc.edu/blazes-of-light-reveal-how-plants-signal-danger-long-distances/}{University of Wisconsin News summary}). GABA also accumulates in plants under many forms of stress, and can directly regulate ion channels in ways that are at least suggestively analogous to its inhibitory roles in animals (\href{https://www.tandfonline.com/doi/full/10.1080/15592324.2020.1862565}{Taylor \& Francis review}; \href{https://www.nature.com/articles/ncomms8879}{Nature Communications: Ramesh et al.~2015}). Dopamine, serotonin, and acetylcholine likewise turn up in plants and algae, sometimes at concentrations that would sound almost comical if one had already decided that these were essentially ``brain chemicals'' (\href{https://pubs.acs.org/doi/10.1021/jf9909860}{ACS Publications: Kanazawa \& Sakakibara 2000}; \href{https://pubmed.ncbi.nlm.nih.gov/10725161/}{PubMed}; \href{https://en.wikipedia.org/wiki/Dopamine}{Wikipedia: Dopamine}; \href{https://auctoresonline.org/article/a-review-on-neurotransmitters-dopamine--serotonin}{Auctores review}; \href{https://pubmed.ncbi.nlm.nih.gov/11243568/}{PubMed: acetylcholine in plants}).

I do not want to overstate the analogy here either. A wounded plant is not having thoughts about its predicament, and a calcium wave in phloem is not a hidden little brain. But these findings are still useful because they undercut the temptation to treat neurotransmitters as substances whose significance begins only once neurons arrive. Caspi et al.~(2023) make that picture much clearer: many neurotransmitters are cheap metabolic intermediates, or close relatives of them, synthesized in very few steps from central metabolic pathways (\href{https://elifesciences.org/articles/83361}{eLife: Caspi et al.~2023}). Glutamate, GABA, dopamine, and acetylcholine all make sense in that light not as magical neural inventions but as ancient, chemically convenient molecules that evolution kept finding new uses for (\href{https://www.creative-proteomics.com/resource/amino-acid-degradation-alpha-ketoglutarate.htm}{Creative Proteomics overview of α-ketoglutarate metabolism}; \href{https://www.tandfonline.com/doi/full/10.1080/15592324.2020.1862565}{Taylor \& Francis review}; \href{https://www.ncbi.nlm.nih.gov/books/NBK20199/}{NCBI Bookshelf overview of catecholamine synthesis}; \href{https://en.wikipedia.org/wiki/Dopamine}{Wikipedia: Dopamine}).

This is one of those points that feels minor until it suddenly does not. If molecules later used in synapses were already sitting deep inside metabolism, stress signaling, and environmental response long before nervous systems existed, then the history of minds begins to look less like the arrival of an enchanted new essence and more like the gradual co-option and recombination of older machinery. Neurons, in that picture, do not invent significance from nothing. They inherit and repurpose it.

That is the main reason I am dwelling on this. The point is not to smuggle the dopamine section forward again. It is to make it harder to keep talking as though the chemistry of mind descended from nowhere, fully formed and morally legible. The chemistry came first. The wiring came later. And once that becomes easier to see, a great many supernaturalized intuitions about mental life start looking less profound and more like leftovers.

\subsection{Phase 3: The deep timeline before anything like a mind existed}\label{03_it_started-phase-3-the-deep-timeline-before-anything-like-a-mind-existed}

This is the part where the chronology gets so large that it risks sounding melodramatic. Unfortunately, the chronology is not my fault.

If I want to tell the story of nervous systems without treating them as a supernatural rupture, then I have to start absurdly early --- before neurons, before animals, arguably before anything we would be tempted to call cognition. That does not mean I think a bacterium has a philosophy of life, or that every feedback loop deserves to be called thought. It means only that many of the operations we later associate with minds --- sensing, comparing, integrating, remembering, acting under uncertainty --- did not appear all at once in a flash of sacred interiority. They were assembled gradually, out of older and simpler machinery.

I am aware that there is a risk here of sounding as though I am redefining ``agency'' until it includes everything from a bacterium to a parliament. I am not trying to do that. I am trying to trace a family resemblance: not consciousness first, then mechanism, but mechanism first, then ever-richer forms of coordination, memory, prediction, and control.

\subsubsection{3.1 Chemotaxis and the temptation to call very old machinery ``agency''}\label{03_it_started-chemotaxis-and-the-temptation-to-call-very-old-machinery-agency}

One of the cleanest examples of pre-neural decision-making is bacterial chemotaxis. \emph{E. coli} cannot sense a spatial gradient across its tiny body, so it does something cleverer: it samples over time. It alternates between runs and tumbles, extending runs when the chemical environment appears to be improving and interrupting them when it does not (\href{https://www.sciencedirect.com/science/article/abs/pii/S0955067411001542}{ScienceDirect: Sourjik \& Wingreen 2012}; \href{https://journals.plos.org/ploscompbiol/article?id=10.1371/journal.pcbi.1005966}{PLOS Computational Biology}). This is often described mathematically as a biased random walk, and the resemblance to stochastic optimization procedures is strong enough that modern readers almost cannot help noticing it (\href{https://arxiv.org/abs/2511.18745}{arXiv}).

I want to be a little careful with that resemblance. Saying that chemotaxis is ``like'' gradient ascent can be illuminating; saying that bacteria are secretly doing machine learning would be embarrassing. Still, the comparison is not empty. The cell senses, compares, updates, and acts. It also implements amplification, receptor integration, and adaptive resetting through methylation and demethylation so that it can remain sensitive across different background concentrations (\href{https://link.springer.com/article/10.1007/s11538-008-9321-6}{Springer overview}). That is already a striking amount of organized problem-solving for something with no neurons, no synapses, and certainly no ghost in the machine.

Whether one wants to call that ``agency'' is partly a matter of taste. I am less attached to the word than to the underlying point: some of the basic operations later folded into nervous systems are far older than nervous systems themselves. Whatever else minds are, they did not begin from nothing.

\subsubsection{3.2 The pharmacy before the wiring: sponges and chemical coordination}\label{03_it_started-the-pharmacy-before-the-wiring-sponges-and-chemical-coordination}

Sponges are useful here because they are so awkward for any story that wants neurons to appear as a clean beginning. They have no neurons, no synapses, and no muscles in the usual sense. And yet they produce and respond to molecules such as glutamate, GABA, nitric oxide, and ATP, and use them in coordinated ways (\href{https://journals.biologists.com/jeb/article/228/3/JEB248010/365649/ATP-and-glutamate-coordinate-contractions-in-the}{Journal of Experimental Biology}; \href{https://journals.biologists.com/jeb/article/213/13/2310/34070/Evidence-for-glutamate-GABA-and-NO-in-coordinating}{Journal of Experimental Biology: Elliott and Leys 2010}). In \emph{Ephydatia muelleri}, for example, bath application of glutamate can trigger coordinated contractions, while GABA plays an inhibitory role that is at least suggestively continuous with later nervous-system functions (\href{https://journals.biologists.com/jeb/article/213/13/2310/34070/Evidence-for-glutamate-GABA-and-NO-in-coordinating}{Journal of Experimental Biology: Elliott and Leys 2010}).

Genomics makes this even harder to ignore. The \emph{Amphimedon queenslandica} genome contains a startling number of orthologues of genes associated with synaptic machinery, suggesting that much of the molecular toolkit later used in nervous systems was already lying around before true neurons emerged (\href{https://www.nature.com/articles/s41598-019-51282-x}{Nature-linked article}). More recently, Musser et al.~identified sponge cell types with secretory vesicles, long processes, and gene-expression profiles suggestive of proto-synaptic organization --- not yet full neurons, but not simply inert blobs of tissue either (\href{https://www.embl.org/ells/2023/11/27/making-sense-of-our-brain-in-the-light-of-evolution/}{EMBL overview}; \href{https://www.researchgate.net/publication/350119637_Encoding_innate_ability_through_a_genomic_bottleneck}{ResearchGate link from the draft}).

This does not prove a neat one-line ancestry from sponge to human synapse. It does, however, make a certain picture easier to defend: the chemistry and much of the cellular apparatus came first, and the tightly wired signaling architecture came later. That is one reason the phrase ``the pharmacy came before the wiring'' continues to feel worth keeping.

\subsubsection{3.3 Receptors as ancient chemical noses}\label{03_it_started-receptors-as-ancient-chemical-noses}

The ancestry of neurotransmitter receptors tells a similar story. G-protein coupled receptors are the largest receptor superfamily in vertebrates, and many metabotropic neurotransmitter receptors use the same basic signal-transduction logic as receptors involved in chemosensation: ligand binding outside, conformational change, G-protein activation, downstream intracellular cascade (\href{https://pmc.ncbi.nlm.nih.gov/articles/PMC4610211/}{PubMed Central review}). In other words, one of the molecular strategies used to ``smell'' the environment was later miniaturized and redirected inward across synaptic clefts.

That is not just a poetic analogy. Sponges and other pre-bilaterian animals already possess members of receptor families that later become central to neurotransmission. Krishnan et al.~found large repertoires of GPCRs in \emph{A. queenslandica}, including homologs related to glutamatergic and GABAergic signaling (\href{https://pmc.ncbi.nlm.nih.gov/articles/PMC10570526/}{PubMed Central}; \href{https://pubmed.ncbi.nlm.nih.gov/25528161/}{PubMed}). Related work showed that receptor families strongly associated with neurotransmission in higher animals are already present in sponges and placozoans (\href{https://pubmed.ncbi.nlm.nih.gov/25696819/}{PubMed}; \href{https://www.researchgate.net/publication/334232136_Sponge_Behavior_and_the_Chemical_Basis_of_Responses_A_Post-Genomic_View}{ResearchGate link from the draft}). Even ionotropic glutamate-receptor relatives show up in protostome olfaction and taste as environmental chemical sensors (\href{https://journals.plos.org/plosgenetics/article?id=10.1371/journal.pgen.1001064}{PLOS Genetics: Croset et al.~2010}).

The upshot, at least as I read it, is not that synapses are trivial. It is that synapses did not invent chemical sensitivity from scratch. They privatized and specialized a much older way of dealing with the world.

\subsubsection{3.4 From secretory cells to neurons}\label{03_it_started-from-secretory-cells-to-neurons}

If one asks where neurons themselves came from, a recurring answer in the literature is: from secretory cells that gradually became more specialized, more excitable, and more tightly connected. Different groups tell this story with different emphases, and there is still room for disagreement about details, including whether neurons evolved once or multiple times. But the broad picture is increasingly hard to dismiss.

SNARE proteins, for instance, are not inventions of synapses. They are part of the general eukaryotic machinery for vesicle fusion and trafficking. The same goes for related components such as Munc18. Choanoflagellates already possess important pieces of this toolkit, and some of the interaction logic is conserved with metazoan synaptic machinery (\href{https://www.ncbi.nlm.nih.gov/pmc/articles/PMC4137706/}{NCBI / PMC}). Voltage-gated calcium and sodium channel homologs also predate neurons, which matters because it means electrical excitability did not suddenly appear as a bonus feature once nervous systems were already in place.

Moroz has argued for a neuronal polygeny hypothesis, according to which neuron-like cell types may have emerged multiple times from different secretory precursors (\href{https://www.frontiersin.org/journals/cell-and-developmental-biology/articles/10.3389/fcell.2021.669087/full}{Frontiers}; \href{https://pmc.ncbi.nlm.nih.gov/articles/PMC8293673/}{PubMed Central}). Jékely, by contrast, emphasizes chemically organized peptidergic signaling networks as a precursor state from which true synapses later emerged (\href{https://pmc.ncbi.nlm.nih.gov/articles/PMC9233960/}{PubMed Central}). Musser's sponge work again provides a tantalizing intermediate: cells with secretory vesicles, projections, and partial pre/post-synaptic gene-expression signatures, suggestive of something like a proto-synapse (\href{https://www.embl.org/ells/2023/11/27/making-sense-of-our-brain-in-the-light-of-evolution/}{EMBL overview}).

The safest conclusion is not that the whole story has been cleanly solved. It is that neurons look less and less like a miraculous invention and more and more like a specialized elaboration of older secretory and excitable systems.

\subsubsection{3.5 Choanoflagellates and the scaffold before the synapse}\label{03_it_started-choanoflagellates-and-the-scaffold-before-the-synapse}

Choanoflagellates are especially annoying to simple origin stories because they keep turning out to have molecular components that sound uncannily ``too advanced'' for single-celled relatives of animals. Proteins associated with the postsynaptic scaffold in animals --- such as Homer, Shank, PSD-95, CaMKII, and others --- have deep evolutionary roots, and choanoflagellate homologs preserve some surprisingly relevant structural features (\href{https://www.nature.com/articles/s41583-025-00983-6}{Nature review}; \href{https://www.researchgate.net/publication/272496772_The_origin_and_evolution_of_synaptic_proteins_-_Choanoflagellates_lead_the_way}{ResearchGate summary of Burkhardt et al.}; \href{https://pmc.ncbi.nlm.nih.gov/articles/PMC2824662/}{PubMed Central}).

That does not mean choanoflagellates have hidden synapses. It means the scaffold later recruited into synapses was already being used for other purposes, especially calcium signaling and cytoskeletal organization. Alié and Manuel's reconstruction of a proto-postsynaptic density captures this nicely: before it became a specialized synaptic structure, it seems to have functioned in more general intracellular coordination (\href{https://pmc.ncbi.nlm.nih.gov/articles/PMC2824662/}{PubMed Central}).

Even cellular memory has ancient roots. The mention in the draft of complex DNA methylation systems in unicellular eukaryotes is relevant here, though I would probably keep the tone modest and avoid sounding as though methylation has solved philosophy. What it does suggest is that cells had ways of preserving state, history, and responsiveness long before neurons appeared (\href{https://phys.org/news/2025-11-celled-complex-dna-epigenetic-code.html}{Phys.org summary of a 2025 \emph{Nature Genetics} result}). Again, not thought, but not nothing.

\subsubsection{3.6 The axon as a solution to the scaling problem}\label{03_it_started-the-axon-as-a-solution-to-the-scaling-problem}

One of the cleanest places where biology starts looking like engineering is the problem of distance. Diffusion is fine across microscopic spaces and terrible across macroscopic ones. A neurotransmitter can cross a synaptic cleft quickly, but diffusion over millimeters or centimeters becomes hopelessly slow because the relevant time scale grows with the square of distance (\href{https://www.physiologyweb.com/calculators/diffusion_time_calculator.html}{PhysiologyWeb diffusion calculator}; \href{https://www1.pbrc.hawaii.edu/~danh/MyelinEvolution/rapid-conduction.html}{Hawaii overview of rapid conduction}).

This makes the emergence of long cellular processes and fast electrical signaling look less like gratuitous biological ornament and more like a practical solution to a brutal scaling problem. Voltage-gated ion channels provided the basis for that solution. Work by Zakon and others suggests that sodium channels evolved from calcium channels, and that important channel families predate true nervous systems (\href{https://www1.pbrc.hawaii.edu/~danh/MyelinEvolution/rapid-conduction.html}{Hawaii overview of rapid conduction}). Once true sodium-selective channels and fast action potentials were available, multicellular organisms could coordinate activity across much larger bodies than diffusion alone would ever allow.

The later evolution of myelination sharpened the same logic rather than changing it. Myelinated vertebrate axons, and analogous solutions in some invertebrates, make clear that whenever evolution runs into the speed limits of passive chemistry, it starts looking for architectural hacks ({[}Hartline \& Colman 2007, as cited in the draft{]}; {[}Hill et al.~2008, as cited in the draft{]}). I would probably soften the phrase ``physics dictates architecture,'' but only slightly. At least here, the pressure exerted by geometry and diffusion really is severe enough to shape what kinds of systems are even viable.

\subsubsection{3.7 Turing patterns and pre-neural computation in body-building}\label{03_it_started-turing-patterns-and-pre-neural-computation-in-body-building}

I was initially unsure whether the Turing-pattern material belonged here, because it risks broadening ``computation'' until the word starts to wobble. But I think it earns its place if handled carefully. Turing's reaction-diffusion model showed that pattern can arise from local chemical interactions without any overseer, through the interplay of activation and inhibition across space. Later work in developmental biology confirmed that this kind of mechanism is not just mathematically cute; it really does help explain features such as pigment patterning, digit formation, hair-follicle spacing, and related developmental regularities (examples in the draft include Watanabe and Kondo 2015; Raspopovic et al.~2014; Sick et al.~2006; Economou et al.~2012).

The reason I think this matters is not that a zebra stripe is thinking. It is that biological form itself already depends on distributed local rules generating global order. By the time nervous systems arrive, biology has long experience building coordinated structure out of decentralized interactions. That does not make nervous systems less special. It just means they did not emerge into a world where coordination had to be invented from scratch.

\subsubsection{3.8 The urbilaterian pivot: from diffuse coordination to centralized prediction}\label{03_it_started-the-urbilaterian-pivot-from-diffuse-coordination-to-centralized-prediction}

At some point, though, something genuinely new does happen. Bilaterian nervous systems are not just bigger sponges or faster cnidarians. The usual picture of the urbilaterian ancestor is that it already possessed a degree of centralization, patterned body axes, and regionally organized neural tissue that distinguish it sharply from diffuse nerve nets (\href{https://pmc.ncbi.nlm.nih.gov/articles/PMC2614231/}{PubMed Central on dorsoventral patterning}; Hirth et al.~2003 as cited in the draft).

This is where the story begins to look less like generic coordination and more like something on the road to prediction, action selection, and model-based control. The lamprey literature is especially important because it suggests that some circuitry associated with vertebrate action selection has very deep roots. Grillner and colleagues argue that lampreys already possess basal ganglia circuits with a striking degree of homology to those of later vertebrates, including dopaminergic modulation ({[}Stephenson-Jones et al.~2011; Grillner \& Robertson 2013; Robertson et al.~2012, all as cited in the draft{]}).

This is exactly the point where I need to keep myself from letting the next section invade this one. The draft is right that these findings matter enormously. It is also right that similar dopaminergic functions show up across bilaterians, from worms to flies to planarians (\href{https://pubmed.ncbi.nlm.nih.gov/17565705/}{PubMed: Nishimura et al.~2007}, plus the additional references cited in the draft). But for the purposes of \emph{this} section, the key point is broader than dopamine: by the time we get to bilaterians and especially early vertebrates, nervous systems are no longer merely transmitting perturbations. They are selecting, modulating, and organizing behavior through centralized architectures that begin to look recognizably like the ancestors of later minds.

That is enough for now. The neurotransmitter-specific consequences can wait a little longer.

\subsubsection{3.9 Dendrites and the embarrassment of the ``simple neuron'' story}\label{03_it_started-dendrites-and-the-embarrassment-of-the-simple-neuron-story}

Another place where modern computational language becomes tempting is the single neuron itself. For a long time, the simplified picture of neurons as near-trivial integrate-and-fire units was useful, and in some contexts still is. But it is also badly incomplete. Dendrites carry out local nonlinear operations, temporal integration, coincidence detection, and sequence sensitivity; they are not just passive wiring for somatic decisions ({[}London and Häusser 2005; Branco, Clark \& Häusser 2010; Poirazi, Brannon, and Mel 2003; Beniaguev et al.~2021, all as cited in the draft{]}).

Calling a pyramidal neuron a ``deep neural network'' is, again, one of those comparisons that can either illuminate or embarrass depending on how triumphantly it is stated. I think it is useful if taken as a reminder that biological computation is richly structured even below the level at which we usually draw our boxes. A single neuron receiving ten thousand or more inputs is not just summing numbers in one place. It is a branched, compartmentalized device with semi-independent local processing. That matters because it undermines the lazy intuition that brains only become interesting once many stupid little units are stacked together. Sometimes one ``unit'' is already hiding a startling amount of architecture.

\subsubsection{3.10 Innateness, pre-training, and evolutionary priors}\label{03_it_started-innateness-pre-training-and-evolutionary-priors}

The final subsection in this phase is probably the most obviously modern in its vocabulary, and therefore the one most in need of self-restraint. Zador's critique of ``pure learning'' makes an important point: organisms do not begin as blank optimization surfaces. The genome cannot specify every synapse individually, so what it gives instead are developmental rules, biases, architectures, and built-in expectations --- compressed priors, if one wants to use that language (\href{https://www.semanticscholar.org/paper/A-critique-of-pure-learning-and-what-artificial-can-Zador/32abbcb3aa75ac34a92624dc779a9f7a82ee981c}{Semantic Scholar link}; \href{https://www.researchgate.net/publication/335310161_A_critique_of_pure_learning_and_what_artificial_neural_networks_can_le}{ResearchGate link from the draft}; \href{https://www.nature.com/articles/s41467-019-11786-6}{Nature-linked genomic bottleneck paper}; \href{https://pmc.ncbi.nlm.nih.gov/articles/PMC11420173/}{PubMed Central: Shuvaev et al.~2024}).

This is one reason I think ML analogies can still be useful, provided they are worn lightly. Retinal waves, for example, really are internally generated training-like signals that help structure later visual function before normal visual experience arrives. That is not merely metaphorical; it is development using endogenous patterned activity to set up future competence ({[}Ackman, Burbridge, and Crair 2012, as cited in the draft{]}; \href{https://arxiv.org/abs/2311.17232}{arXiv on retinal-wave pretraining}; \href{https://www.sciencedirect.com/science/article/abs/pii/S0010027725003567?dgcid=rss_sd_all}{ScienceDirect: Wood and Wood 2025}). Likewise, innate odor responses in \emph{Drosophila} show that some behavioral valences are already heavily scaffolded by circuit design rather than learned from scratch (\href{https://www.cell.com/fulltext/S0092-8674(12)01357-8}{Cell: Stensmyr et al.~2012}; \href{https://link.springer.com/article/10.1007/s00441-020-03392-6}{Springer overview}).

Friston's language of priors and updating belongs here too, though again I would rather use it as a suggestive framework than as scripture (\href{https://www.semanticscholar.org/paper/The-free-energy-principle:-a-unified-brain-theory-Friston/1ed6a4a10589618d4f26350f1a296ee767ceff6b}{Semantic Scholar: Free Energy Principle}). The safe and important point is that organisms inherit structured expectations from evolution and then modify them through experience. They do not begin nowhere.

And that, really, is the deeper throughline of this entire phase. The machinery of mind was not dropped into the world fully formed. It was layered together across absurd spans of time: chemical sensitivity, secretory coordination, excitability, cellular memory, long-distance transmission, centralization, branched local computation, developmental priors, and only then the increasingly sophisticated forms of learning and behavioral control that later chapters will care about more directly. I am not claiming that this sequence dissolves every mystery. I am claiming that it makes magical thinking about the mind harder to defend.

Which is enough. For now.

\subsubsection{Bridge to the next section}\label{03_it_started-bridge-to-the-next-section}

The deeper throughline of this section is not that every living thing has a mind, still less that modern machine learning has somehow solved biology in reverse. It is that the machinery from which minds were eventually built did not appear all at once. Many of its supposedly ``special'' components have older, non-neural histories: signaling molecules with deep metabolic roots, receptor systems repurposed from environmental sensing, secretory machinery older than synapses, electrical excitability older than nervous systems, and centralized architectures assembled gradually rather than dropped into the world whole. The chemistry came first. The wiring came later. And once that becomes easier to see, magical thinking about the mind becomes harder to defend.

None of this dissolves every mystery, and I am not pretending it does. But it does make one thing much less plausible: the idea that failures of motivation, action, or behavioral control are best understood as moral stains on an otherwise immaterial will. If minds are built, then they can break. If agency is implemented, then it can be impaired. That is enough to justify the next question, which is not whether one particular neurotransmitter explains the whole of human life, but whether some parts of this machinery matter so much that damaging them predictably dismantles the organism's ability to pursue its own ends.

\subsection{Opening move for the next section}\label{03_it_started-opening-move-for-the-next-section}

What follows next belongs to a different level of argument. Up to this point, I have been trying to show that minds have a material history, and that the concepts we use to describe mental life should reflect that history rather than smuggling magic back in. The next section narrows the focus. It asks whether one particular part of this inherited machinery matters so much that impairing it produces a recognizable collapse in initiative, reward-guided behavior, effort allocation, and the organism's ability to orient itself toward its own future.

That is where the dopamine argument properly begins. Not because dopamine explains everything, and not because I want to turn one neurotransmitter into a secular soul, but because there are cases in which the folk-psychological picture breaks down so cleanly that continuing to talk in terms of character or will starts to look less like common sense and more like superstition with good manners. If depriving a human being of the relevant machinery predictably dismantles their ability to act, then what is being damaged is not some decorative extra. It is a substantial part of the substrate of agency itself.
